\documentclass[12pt,a4paper,UTF8]{ctexart}
\usepackage[a4paper,top=22mm,bottom=22mm,left=22mm,right=22mm,headheight=15pt]{geometry}
\usepackage{amsmath,amssymb,mathtools,bm}
\usepackage{graphicx,xcolor,booktabs,longtable,array,enumitem,fancyhdr,caption,float,microtype,setspace}
\newcommand{\Needspace}[1]{\par\begingroup\dimen0=#1\relax\dimen2=\pagegoal\advance\dimen2 by-\pagetotal\ifdim\dimen0>\dimen2\newpage\fi\endgroup}
\newcommand{\needspace}[1]{\par\begingroup\skip0=\lastskip\vskip-\skip0\dimen0=#1\relax\advance\dimen0 by\skip0\vskip\dimen0\penalty-100\vskip-\dimen0\vskip\skip0\penalty9999\endgroup}
\usepackage[hidelinks,unicode]{hyperref}
\usepackage[most]{tcolorbox}
\definecolor{ContestBlue}{HTML}{17365D}
\hypersetup{pdftitle={第042套},pdfauthor={竞赛物理教研}}
\setstretch{1.24}
\setlength{\parindent}{2em}
\setlength{\parskip}{0.35em}
\setlength{\emergencystretch}{2em}
\setlist{itemsep=0.35em,topsep=0.4em}
\captionsetup{font=small,labelfont=bf,justification=centering}
\setcounter{secnumdepth}{0}
\pagestyle{fancy}
\fancyhf{}
\fancyhead[L]{\small 第042套}
\fancyfoot[C]{\thepage}
\renewcommand{\headrulewidth}{0.4pt}
\ctexset{
  section={format=\Large\bfseries\color{ContestBlue},beforeskip=1.2em,afterskip=0.7em},
  subsection={format=\large\bfseries\color{ContestBlue},beforeskip=1.0em,afterskip=0.5em},
  subsubsection={format=\normalsize\bfseries,beforeskip=0.8em,afterskip=0.35em}
}
\makeatletter
\newcommand{\tightlist}{\setlength{\itemsep}{0pt}\setlength{\parskip}{0pt}}
\newcommand{\maxwidth}{\ifdim\Gin@nat@width>\linewidth\linewidth\else\Gin@nat@width\fi}
\newcommand{\maxheight}{\ifdim\Gin@nat@height>0.72\textheight 0.72\textheight\else\Gin@nat@height\fi}
\newcommand{\pandocbounded}[1]{\begingroup\setkeys{Gin}{width=\maxwidth,height=\maxheight,keepaspectratio}#1\endgroup}
\makeatother
\begin{document}
\begin{titlepage}
  \centering
  \vspace*{58mm}
  {\Huge\bfseries\color{ContestBlue} 第042套\par}
  \vfill
  \vspace*{28mm}
\end{titlepage}
\section{第 1 题}\label{ux7b2c-1-ux9898}

孙鹏的体重超过体重秤量程。他找来两块与秤等高的砖和一块钢板，把三个支承排成等边三角形，站在钢板上标为
P 的位置，读数为 \(G_1\)；原卷写``轮换体重秤和钢板''，又得到
\(G_2,G_3\)。判断下列说法中哪些不正确：

A．只有 P 在三角形中心，\(G_1+G_2+G_3\) 才等于人和钢板的总重；

B．只有支点构成等边三角形，三读数之和才等于总重；

C．各读数与对应顶点到 P 的距离成反比；

D．以 P 为原点，三支点位置向量为 \(\boldsymbol r_i\)，有
\(\sum_iG_i\boldsymbol r_i=0\)。

题面中钢板自重、P
与合成重心的关系及轮换操作存在歧义，解答分别说明通常的固定载荷模型和字面模型。

\begin{figure}[H]
\centering
\includegraphics[width=0.9\linewidth,height=\textheight,keepaspectratio,alt={题面模型}]{../figures/question-01.pdf}
\caption{题面模型}
\end{figure}

\newpage

\section{第 2 题}\label{ux7b2c-2-ux9898}

双黑洞系统由 \(29M_\odot\) 和 \(36M_\odot\)
的黑洞组成，\(M_\odot=1.99\times10^{30}\,\mathrm{kg}\)。假设作牛顿圆轨道运动，轨道频率
\(30\,\mathrm{Hz}\)，引力波频率为轨道频率的两倍，(1)
求两黑洞中心距离。(2) 按原题给定的半径
\(r_i=Gm_i/c^2\)，估计两``表面''接触时的轨道频率。取
\(G=6.67\times10^{-11}\)（SI）、\(c=3.0\times10^8\,\mathrm{m/s}\)。

原卷将 \(Gm/c^2\) 称为史瓦西半径；标准非旋转黑洞史瓦西半径为
\(2Gm/c^2\)。解答先按题设半径估算，再列出因子 2 修正的结果。

\begin{figure}[H]
\centering
\includegraphics[width=0.9\linewidth,height=\textheight,keepaspectratio,alt={题面模型}]{../figures/question-02.pdf}
\caption{题面模型}
\end{figure}

\newpage

\section{第 3 题}\label{ux7b2c-3-ux9898}

在海水深 \(H=2\,\mathrm{km}\) 处产生一个初始体积 \(V_0=1\,\mathrm L\)
的气泡，并同时释放一个密封刚性空瓶，瓶中空气的质量和瓶的排水体积分别与初始气泡相同。海面气压
\(p_0=10^5\,\mathrm{Pa}\)，海面空气密度
\(\rho_a=1.21\,\mathrm{kg/m^3}\)。气体等温膨胀，忽略黏滞阻力，求：(1)
到海面时气泡体积；(2) 深度 \(h\) 处气泡和钢瓶各自获得的净能量；(3)
二者到海面时的速度。

海水密度和 \(g\) 未在题面中给出，符号结果保留 \(\rho,g\)。仅为数值示例取
\(\rho=1000\,\mathrm{kg/m^3},g=10\,\mathrm{m/s^2}\)。瓶的质量按题设总质量等于气泡质量理解，不另加未给出的瓶壁质量。

\begin{figure}[H]
\centering
\includegraphics[width=0.9\linewidth,height=\textheight,keepaspectratio,alt={题面模型}]{../figures/question-03.pdf}
\caption{题面模型}
\end{figure}

\newpage

\section{第 4 题}\label{ux7b2c-4-ux9898}

波长为 \(\lambda_0\) 的单色光穿过双折射薄片，两本征线偏振在片内波长为
\(\lambda_1,\lambda_2\)，折射率为
\(n_1,n_2\)。求四分之一波片的最小厚度，并用
\(\lambda_0=589\,\mathrm{nm},n_1=1.875,n_2=1.635\)
求数值。讨论两分量的相位差和入射偏振方向对出射偏振的影响。

\begin{figure}[H]
\centering
\includegraphics[width=0.9\linewidth,height=\textheight,keepaspectratio,alt={题面模型}]{../figures/question-04.pdf}
\caption{题面模型}
\end{figure}

\newpage

\section{第 5 题}\label{ux7b2c-5-ux9898}

一光滑轨道通过原点，近原点形状为
\(y=x^2/a\)，质点在重力作用下作小振动，求周期
\(T_a\)。若要求从任意允许高度由静止释放都有相同周期
\(T_a\)，求轨道高度与弧长的关系，并给出轨道参数方程。

将质点换为半径 \(r\)
的均匀圆柱，要求在足够静摩擦下纯滚动且各允许释放幅度的周期仍为
\(T_a\)，求相应接触轨道的参数方程。

\begin{figure}[H]
\centering
\includegraphics[width=0.9\linewidth,height=\textheight,keepaspectratio,alt={题面模型}]{../figures/question-05.pdf}
\caption{题面模型}
\end{figure}

\newpage

\section{第 6 题}\label{ux7b2c-6-ux9898}

宇宙飞船初始静质量为 \(m_0\)，在实验室中以恒定速率 \(v\)
运动。用燃料反应后喷出的物质改变运动方向，喷出物相对飞船的速度大小为
\(u\)，不作 \(u\ll c\) 或 \(v\ll c\) 近似。燃料反应前的微小静质量为
\(dm_0\)。求最优喷射方向下偏转角 \(d\theta\) 与 \(dm_0\)
的关系；若保持速率转过角度 \(\theta\)，末静质量为 \(m_1\)，求最小质量比
\(m_0/m_1\)。

\begin{figure}[H]
\centering
\includegraphics[width=0.9\linewidth,height=\textheight,keepaspectratio,alt={题面模型}]{../figures/question-06.pdf}
\caption{题面模型}
\end{figure}

\newpage

\section{第 7 题}\label{ux7b2c-7-ux9898}

质量为 \(m\)、电荷为 \(q\)
的小球位于水平硬壁上方，均匀电场使其受到恒定向下的力
\(qE\)，不计重力，与壁的碰撞完全弹性。小球从高度 \(h\)
静止释放后反复弹跳。电场强度缓慢增大到原来的两倍，求新最高点高度与弹跳周期，并讨论``缓慢''的条件。

\begin{figure}[H]
\centering
\includegraphics[width=0.9\linewidth,height=\textheight,keepaspectratio,alt={题面模型}]{../figures/question-07.pdf}
\caption{题面模型}
\end{figure}

\newpage

\section{第 8 题}\label{ux7b2c-8-ux9898}

一均匀实心球质量为 \(m\)、半径为 \(r\)，球面上有一质量也为 \(m\)
的点状兔子。初始系统无角动量，球的体坐标轴与实验室坐标轴重合；兔子位于
\((r\sin\theta,0,r\cos\theta)\)，\(0<\theta<90^\circ\)。

\begin{enumerate}
\def\labelenumi{(\arabic{enumi})}
\item
  球心固定，兔子在实验室中沿高度为 \(r\cos\theta\)
  的水平圆周跑一圈，哪些 \(\theta\) 使球的体轴 \(z'\) 回到初始方向？
\item
  球心可自由运动且系统初始质心静止，兔子仍在实验室沿高度为
  \(r\cos\theta\) 的圆周跑一圈，同问。
\item
  一般 \(\theta\) 下，求前两情形的一圈后 \(z'\) 方向的极角和方位角。
\item
  球心固定，兔子改为相对于球体沿固定纬度圆跑一圈，描述实验室中球轴
  \(z'\) 和兔子的运动。
\end{enumerate}

\begin{figure}[H]
\centering
\includegraphics[width=0.9\linewidth,height=\textheight,keepaspectratio,alt={题面模型}]{../figures/question-08.pdf}
\caption{题面模型}
\end{figure}
\end{document}
